\clearpage
\ExposeTOC
\begin{center}
{\large SGA 7 --- 第 XVIII 讲}\par
\vspace{2em}
{\Large\underline{Lefschetz 线束的上同调研究}}\par
\vspace{1.5em}
{\large N. Katz\footnote{根据 Grothendieck 的简要笔记写成.}}
\end{center}

\section*{目录}

\noindent 1. 射影丛的上同调. \hfill 2

\noindent 2. 爆破簇的上同调. \hfill 4

\noindent 3. 与线束相伴的爆破. \hfill 13

\noindent 4. 与线束相伴的爆破的上同调. \hfill 17

\noindent 5. 与线束相伴的态射 $\rho:\widetilde X\longrightarrow\mathbb P^1$ 的上同调研究.
\hfill 24

\noindent\quad 5.1. 限制同态
$H^*(\widetilde X)\longrightarrow H^*(Y)$. \hfill 24

\noindent\quad 5.2. “强 Lefschetz 定理” (Lefschetz vache, LV). \hfill 27

\noindent\quad 5.3. 条件 (A). \hfill 33

\noindent\quad 5.8. Griffiths 同态. \hfill 47

\noindent\quad 5.9. 不满足假设 (A) 的一个例子.
\hfill 57

\noindent 6. 消灭圈局部理论的应用
(Picard--Lefschetz). \hfill 58

\section*{引言}

本讲给出 Lefschetz 线束的初等上同调结构. 第 1--4 节给出有关射影丛和爆破的上同调结构的“熟知”事实. 第 5 节在假设 $X$ 满足“强 Lefschetz 定理” (5.2), 且所考虑的线束 $D$ 满足某个在 $D$ 上具有局部性质的条件 (A) (5.3) 的前提下, 进行标题所述的研究. 主要结果是 5.6、5.7 和 5.8.7; 最后一个结果在第 XX 讲 Griffiths 定理的证明中起着关键作用. 最后一节明显较不初等, 我们利用 SGA 1 中 Mme Raynaud 关于\underline{驯基本群}的第 XIII 讲, 以及\underline{Picard--Lefschetz 公式} (XV 3.4), 证明:
\begin{enumerate}[label=(\arabic*)]
\item 条件 (A) 往往成立 (6.3, 6.4);
\item 除 $\dim(X)$ 为奇数且特征为 $2$ 的情形外, Lefschetz 线束的单值群在“消灭上同调”上的作用不可约, 并在经由 Picard--Lefschetz 公式与线束的各个“临界点”相伴的“局部消灭圈”集合上的作用可迁 (相差符号视为相同) (6.6 和 6.7).
\end{enumerate}

\begin{editorialnote}
原印本及英译底本在引言中将 Picard--Lefschetz 公式指向第 XVI 讲. 该讲讨论 Milnor 公式; 此处所用的是第 XV 讲 3.4 的 Picard--Lefschetz 公式, 故订正引用. 原印本位置为本讲第 1 页 (全书第 255 页).
\end{editorialnote}

\section*{1. 射影丛的上同调}

\subsubsection*{定理 1.1 (Deligne).}
\underline{设 $f:X\longrightarrow S$ 为固有态射, $\ell$ 为在 $S$ 上可逆的素数, $u\in H^2(X,\mathbb Z_\ell(1))$ (分别地, $u\in H^2(X,\mathbb Q_\ell(1))$) 为 $X$ 上的上同调类, $\mathfrak F$ 为 $X$ 上的 $\mathbb Z_\ell$-层 (分别地, $\mathbb Q_\ell$-层). 假设存在整数 $n$, 使得对每个整数 $j$ 和每个整数 $i>0$, 映射}
\begin{equation}\tag{1.1.1}
R^{n-i}f_*\bigl(\mathfrak F(j)\bigr)
\xrightarrow{\ \wedge u^i\ }
R^{n+i}f_*\bigl(\mathfrak F(j+i)\bigr)
\end{equation}
\underline{是 $S$ 上的 $\mathbb Z_\ell$-层 (分别地, $\mathbb Q_\ell$-层) 的同构. 则对每个态射 $g:S\longrightarrow T$, Leray 谱序列}
\[
E_2^{p,q}=R^pg_*R^qf_*(\mathfrak F)
\Longrightarrow R^{p+q}(gf)_*(\mathfrak F)
\]
\underline{退化.}

证明见 [1, 命题~2.4]. 注意, 由固有基变换定理 (SGA 4 XII, 5.2), 可以逐纤维检验此假设.

\subsubsection*{定理 1.2 (参见 SGA 5 VII, 2.2.1).}
\underline{设 $\mathcal E\longrightarrow X$ 为秩为 $1+r$ 的局部自由向量丛. 以 $f:\mathbb P(\mathcal E)\longrightarrow X$ 表示 $\mathcal E$ 定义的射影丛, 并设 $u\in H^2(\mathbb P(\mathcal E),\mathbb Z_\ell(1))$ 为典范可逆层 $\mathcal O_{\mathbb P(\mathcal E)}(1)$ 的第一 Chern 类. 对在 $X$ 上可逆的素数 $\ell$, 有:}
\begin{enumerate}[label=\alph*)]
\item \underline{谱序列}
\begin{equation}\tag{1.2.1}
E_2^{p,q}=H^p\bigl(X,R^qf_*\mathbb Z_\ell\bigr)
\Longrightarrow
H^{p+q}\bigl(\mathbb P(\mathcal E),\mathbb Z_\ell\bigr)
\end{equation}
\underline{退化;}
\item \underline{映射}
\begin{equation}\tag{1.2.2}
\bigoplus_{0\leq j\leq r}
H^{q-2j}\bigl(X,\mathbb Z_\ell(-j)\bigr)
\longrightarrow
H^q\bigl(\mathbb P(\mathcal E),\mathbb Z_\ell\bigr),
\qquad
(\tau_{q-2j})_j\longmapsto
\sum_j f^*(\tau_{q-2j})\wedge u^j
\end{equation}
\underline{是同构.}
\end{enumerate}
\footnote{在 (1.2.2) 中, 印本把上限写成 $n$, 把次数写成 $q-2-j$, 并少了一个右括号. 秩 $1+r$、定义映射以及紧接在下面的证明都要求 $0\leq j\leq r$ 和 $q-2j$; 此处给出改正后的公式.}

\noindent\underline{证明.}
$f$ 的纤维是 $r$ 维射影空间, 而我们知道
\begin{equation}\tag{1.2.3}
R^if_*\mathbb Z_\ell
\xleftarrow[\ \sim\ ]{\ \wedge^{i/2}u\ }
\begin{cases}
\mathbb Z_\ell\left(-\dfrac{i}{2}\right),
  & \text{若 $i$ 为偶数且 $0\leq i\leq2r$},\\[2ex]
0, & \text{其他情形}.
\end{cases}
\end{equation}
将定理 1.1 应用于 $(f,\ell,u,\mathbb Z_\ell,r)$, 即得 (a). 由 (1.2.3),
\begin{equation}\tag{1.2.4}
E_2^{a,b}=
\begin{cases}
H^a\left(X,\mathbb Z_\ell\left(-\dfrac b2\right)\right),
  & \text{若 $b$ 为偶数且 $0\leq b\leq2r$},\\[2ex]
0, & \text{其他情形}.
\end{cases}
\end{equation}
因此, 对谱序列给出的滤过, 与 $H^q(\mathbb P(\mathcal E),\mathbb Z_\ell)$ 相伴的分次对象为
\[
\bigoplus_{0\leq j\leq r}
H^{q-2j}\bigl(X,\mathbb Z_\ell(-j)\bigr).
\]
若在此直和上赋予滤过
\[
F^i=
\bigoplus_{\substack{0\leq j\leq r\\q-2j\leq i}}
H^{q-2j}\bigl(X,\mathbb Z_\ell(-j)\bigr),
\]
则映射 (1.2.2) 与滤过相容, 并在相伴分次对象上诱导恒等映射. 这就证明了 (b).

